\documentclass{article}

%Package-free
\setlength{\unitlength}{1em}
\newcommand\like[1]{\begin{picture}(1,1)
\ifnum0=#1\put(.5,.35){\circle{1}}\else
\ifnum100=#1\put(.5,.35){\circle*{1}}\else
\put(.5,.35){\circle{1}}\put(.5,.35){\circle*{.#1}}
\fi\fi\end{picture}}

\usepackage{tikz}

\newcommand*{\kreise}[1]{\begin{tikzpicture}[scale=0.15]%
    \draw (0,0) circle (1);
    \fill[fill opacity=1,fill=black] (0,0) circle (#1/100);
  \end{tikzpicture}}

\newcommand*{\torte}[1]{\begin{tikzpicture}[scale=0.15]%
    \draw (0,0) circle (1);
    \fill[fill opacity=1,fill=black] (0,0) -- (90:1) arc (90:90-#1*3.6:1) -- cycle;
  \end{tikzpicture}}

\newcommand{\rating}[1]{%
  \begin{tikzpicture}[x=1ex,y=1ex]
    \begin{scope}
      \clip (0,1) circle (1);
      \fill[black] (-1,0) rectangle (1,#1/50);
    \end{scope}
    \draw[black, thin, radius=1] (0,1) circle;
  \end{tikzpicture}%
}

\newcommand{\ratingb}[2][1ex]{%
  \pgfmathsetmacro\th{asin(#2/50-1)}% (theta angle of polar coordinates)
    \tikz{%
      \fill[black] (\th:#1) arc (\th:-180-\th:#1) -- cycle;
      \draw[black, thin, radius=#1] (0,0) circle;
    }%
}

\usetikzlibrary{calc}
\tikzset{circle mark/.pic={
    \draw [path picture={%
      \fill (path picture bounding box.south west) |-  
      ($(path picture bounding box.south east)!#1/100!(path picture bounding box.north east)$) 
      |- cycle;}] circle [radius=1ex];
  }}

\usetikzlibrary{shapes.geometric}
%for triangle

\begin{document}

\foreach \i in {0,5,...,100} {\like{\i}}\\

\foreach \i in {0,5,...,100} {\kreise{\i}}\\

\foreach \i in {0,5,...,100} {\torte{\i}}\\

\foreach \i in {0,5,...,100} {\rating{\i}}\\

\foreach \i in {0,5,...,100} {\ratingb{\i}}\\

\foreach \i in {0,5,...,100}
  \tikz\path pic{circle mark=\i};

\begin{tikzpicture}[scale=0.15]
  \fill[black] (0,0) -- (2,0) -- (1,-2) -- cycle;
\end{tikzpicture}
\begin{tikzpicture}[scale=0.15]
  \draw (0,0) rectangle (2,2);
\end{tikzpicture}

\tikz \fill[black] (1ex,1ex) circle (1ex);

\tikz \draw (0,0) -- (1,0) -- (1,1) -- cycle;

\begin{tikzpicture}[even odd rule,rounded corners=2pt,x=10pt,y=10pt]
\filldraw[fill=black] (0,0) rectangle (1,1)
[xshift=5pt,yshift=5pt] (0,0) rectangle (1,1);
\end{tikzpicture}

\tikz \fill[black,rounded corners=2pt] (0,0) rectangle (1em,1em);

\begin{tikzpicture}
  \draw[rounded corners=2pt] (0,0) rectangle (1em,1em);
  \fill[black] (.5em,.5em) -- (1em,.5em) {[rounded corners=2pt] -- (1em,1em)} -- (.5em,1em) -- cycle;
\end{tikzpicture}

\begin{tikzpicture}
  \draw[rounded corners=2pt] (0,0) rectangle (1em,1em);
  \fill[black, rounded corners=2pt] (.2em,.2em) rectangle (.8em,.8em);
\end{tikzpicture}

\begin{tikzpicture}
  \draw[rounded corners=2pt] (0,0) rectangle (1em,1em);
  \fill[black, radius=.1em] (.5em,.5em) circle;
\end{tikzpicture}

\begin{tikzpicture}[scale=0.15]
  \fill[black, rounded corners=2pt] (0,0) -- (2,0) -- (1,2) -- cycle;
\end{tikzpicture}
\begin{tikzpicture}[scale=0.15]
  \draw[rounded corners=2pt] (0,0) -- (2,0) -- (1,2) -- cycle;
\end{tikzpicture}
\begin{tikzpicture}[scale=0.15]
  \draw[radius=1] (0,0) circle;
\end{tikzpicture}

\begin{tikzpicture}[scale=0.15]
  \fill[black, rounded corners=2pt] (0,0) -- (2,0) -- (1,-2) -- cycle;
\end{tikzpicture}
\begin{tikzpicture}[scale=0.15]
  \fill[black] (0,0) arc (25:-205:1) -- cycle;
\end{tikzpicture}
\begin{tikzpicture}
  \fill[black, rounded corners=1.5pt] (0,0) rectangle (1em,3pt);
\end{tikzpicture}


\end{document}
